\documentclass[12pt]{article}
\usepackage{amsmath, amssymb, graphicx}
\usepackage{geometry}
\geometry{margin=2.5cm}

\title{\textbf{Question 4: Neural Network Derivation}}
\date{}

\begin{document}

\maketitle

\section*{Question 4}

\subsection*{Problem Statement}

Consider the neural network shown in Figure 5 for a binary classification task. Assume the hidden layer uses a linear activation function:
\[ h(z) = cz \]
and the output layer uses the sigmoid function:
\[ g(z) = \frac{1}{1 + e^{-z}} \]

The network wants to learn a function for:
\[
P(Y = 1 \mid X, W)
\quad \text{where} \quad W = (w_1, w_2, \ldots, w_9), \quad X = (x_1, x_2)
\]

\begin{center}
    \includegraphics[width=0.6\textwidth]{1.png} \\
    Figure 5: The neural network architecture for Question 4
\end{center}

\subsection*{1. Derive the final output of the network \( P(Y=1 \mid X, W) \)}

First, compute the activations at the hidden layer:
\begin{align*}
Z_1' &= w_3 x_1 + w_5 x_2 + w_1 \\
Z_2' &= w_4 x_1 + w_6 x_2 + w_2
\end{align*}

Then compute the output before the activation function:
\[
y' = (cw_8 w_3 + cw_9 w_4)x_1 + (cw_8 w_5 + cw_9 w_6)x_2 + cw_1 + cw_2 + w_7
\]

Final output:
\[
y = g(y') = \frac{1}{1 + e^{-y'}}
\]

\subsection*{2. Write the decision boundary}

On the decision boundary:
\[
y = \frac{1}{1 + e^{-y'}} = 0.5 \Rightarrow y' = 0
\]
Thus:
\[
a_1 x_1 + a_2 x_2 + b = 0
\]
where:
\begin{align*}
a_1 &= cw_8 w_3 + cw_9 w_4 \\
a_2 &= cw_8 w_5 + cw_9 w_6 \\
b &= cw_1 + cw_2 + w_7
\end{align*}

\subsection*{3. Can this be done without a hidden layer?}

Yes, if a direct linear combination of \( x_1 \) and \( x_2 \) using appropriate weights could approximate the same decision boundary, it is possible to eliminate the hidden layer. However, the hidden layer provides more flexibility in representing complex boundaries.

\end{document}